\documentclass[a4paper,12pt]{article}
\usepackage[utf8]{inputenc}
\usepackage[T1]{fontenc}
\usepackage[english,russian]{babel}
\usepackage{amsmath, amssymb, amsthm}
\usepackage{geometry}
\geometry{top=2cm, bottom=2cm, left=2.5cm, right=2.5cm}
\usepackage{hyperref}
\hypersetup{colorlinks=true, linkcolor=blue, urlcolor=blue}
\usepackage{indentfirst}
\usepackage{enumitem}

\title{Топологическая связь в Теории Мод:\\ принципы, механизмы, применения}
\author{Теория Мод (автор)}
\date{Версия 2.0 --- 15 июня 2026}

\begin{document}

\maketitle

\begin{abstract}
\noindent
Документ представляет полное описание топологической связи — канала передачи энергии и информации между двумя идентичными топологическими антеннами, который не подчиняется классическим ограничениям скорости света и не зависит от расстояния в эмерджентном пространстве-времени. В рамках Теории Мод этот эффект выводится из прямых корреляций между собственными модами идентичных узлов сети. Документ самодостаточен: все используемые понятия определены заново, без ссылок на внешние тексты. Приведены математическая модель, инженерные требования к антеннам, протоколы модуляции, оценки энергетики и экспериментальные предсказания. Отдельно рассмотрены два режима: передача энергии (требует только кристаллографической идентичности) и передача информации (требует дополнительно идентичности ближнего электромагнитного окружения). Обсуждается связь с квантовой запутанностью и возможное применение для сверхдальней беспроводной передачи энергии.
\end{abstract}

\noindent
\textbf{Ключевые слова:} топологическая связь, Теория Мод, идентичные антенны, топологический изолятор, передача быстрее света, корреляция мод, кристаллографическая идентичность, ближнее окружение.

\tableofcontents

\section*{Часть 1. Фундаментальные понятия (даются заново)}
\addcontentsline{toc}{section}{Часть 1. Фундаментальные понятия (даются заново)}

\subsection{Частица}
Частица — элементарный узел сети. Сеть не имеет координат. Частицы различны, но могут быть идентичны по своим внутренним характеристикам.

\subsection{Мода}
Мода \(M_{A \to B}\) — это упорядоченная структура, связывающая две различные частицы \(A\) и \(B\). Она имеет две компоненты:
\begin{enumerate}
    \item \textbf{Продольная структура} — дискретный путь от \(A\) к \(B\):
    \[
    \Gamma_{AB} = (A = P_0, P_1, P_2, \dots, P_d = B)
    \]
    где \(d \ge 0\) — количество промежуточных частиц-посредников. Если \(d = 0\), частицы находятся в прямом контакте. Если путь не существует, считается \(d = \infty\).
    \item \textbf{Поперечный спектр} — набор вещественных чисел, сопоставленных шагам пути:
    \[
    F_{A\to B}(0), F_{A\to B}(1), \dots, F_{A\to B}(d)
    \]
    Значение \(F_{A\to B}(0)\) называется \textbf{проекцией} \(s_{AB}\). Значения \(F_{A\to B}(k)\) при \(k \ge 1\) описывают распределение взаимодействия вдоль цепочки посредников.
\end{enumerate}

\subsection{Встречные моды и корреляция}
Для каждой пары \((A, B)\), для которой существуют оба пути \(\Gamma_{AB}\) и \(\Gamma_{BA}\), определяются встречные моды. Корреляция вычисляется как свертка поперечных спектров вдоль пути:
\[
C_{AB} = \sum_{k=0}^{d} F_{A\to B}(k) \cdot F_{B\to A}(d-k)
\]
Если \(d = 0\) (прямой контакт), то \(C_{AB} = s_{AB} \cdot s_{BA}\).

Величина \(|C_{AB}|\) называется \textbf{нескомпенсированностью} пары. Чем больше нескомпенсированность, тем сильнее взаимодействие.

\subsection{Собственная мода частицы}
Частица \(A\) имеет \textbf{замкнутую} собственную моду \(M_{A\to A}\). Это путь, начинающийся и заканчивающийся на \(A\):
\[
\Gamma_{AA} = (A = P_0, P_1, \dots, P_{d_{AA}} = A)
\]
Длина пути \(d_{AA} \ge 1\) (минимальная петля — через одну промежуточную частицу). Поперечный спектр собственной моды обозначается \(F_{A\to A}(k)\), \(k = 0, \dots, d_{AA}\), с условием замкнутости \(F_{A\to A}(0) = F_{A\to A}(d_{AA})\).

Собственная мода описывает внутреннюю структуру частицы — её «виртуальное облако» посредников и флуктуации. Важнейшая характеристика собственной моды — \textbf{спиновая последовательность} (знаки \(F_{A\to A}(k)\)), которая определяет внутреннюю ориентацию в спиновом пространстве.

\subsection{Идентичность частиц (топологическая, кристаллографическая, окруженческая)}
Две частицы \(A\) и \(B\) называются \textbf{топологически идентичными}, если:
\begin{enumerate}
    \item Их собственные спектры совпадают как векторы в пространстве всех проекций:
    \[
    s_{A j} = s_{B j} \quad \text{для всех } j \neq A,B
    \]
    (проекции на одинаковые типы частиц равны).
    \item Их собственные моды \(M_{A\to A}\) и \(M_{B\to B}\) имеют одинаковые длины, одинаковые последовательности знаков поперечного спектра и одинаковые топологические инварианты (число намоток, зацеплений и т.д.).
    \item \textbf{Кристаллографическая идентичность}: материал антенны имеет одинаковую кристаллическую структуру и одинаковую ориентацию кристаллографических осей относительно структуры антенны. Это фиксирует спиновую последовательность абсолютно (не относительно эмерджентного пространства).
    \item \textbf{Идентичность ближнего окружения} (требуется только для передачи информации, не для передачи энергии): токи, наведённые в проводниках, близлежащие электрические поля, температуры и другие факторы, влияющие на локальную нескомпенсированность, должны быть одинаковыми для обеих антенн (или, по крайней мере, стационарными и взаимно известными). Это обеспечивает возможность модуляции корреляции без паразитных шумов.
\end{enumerate}

В случае топологических антенн, изготовленных из кристаллических материалов, идентичность достигается производством на одной пластине и размещением в одинаковых электромагнитных условиях (экранирование, одинаковое питание, одинаковая термостабилизация). \textbf{Эмерджентное расстояние и взаимная ориентация в пространстве не входят в условия идентичности.}

\subsection{Корреляция идентичных частиц}
Если \(A\) и \(B\) топологически идентичны (условия 1–3 выше), то их собственные моды совпадают с точностью до циклического сдвига индекса:
\[
F_{A\to A}(k) = F_{B\to B}(k + \delta) \quad \text{(по модулю длины)}
\]
Корреляция между их собственными модами вычисляется как:
\[
C_{AB}^{\text{(собств)}} = \sum_{k=0}^{d_{AA}} F_{A\to A}(k) \cdot F_{B\to B}(d_{AB} - k)
\]
В силу идентичности корреляция \textbf{не зависит от эмерджентного расстояния \(d_{AB}\)} и от эмерджентной ориентации антенн.

Если дополнительно выполнено условие 4 (идентичность ближнего окружения), то корреляция стабильна во времени и может быть модулирована внешним воздействием на передатчик. Если окружение разное, базовая корреляция сохраняется, но её модуляция становится неэффективной из-за несогласованных шумов.

\section*{Часть 2. Эмерджентное пространство-время как приближение}
\addcontentsline{toc}{section}{Часть 2. Эмерджентное пространство-время как приближение}

\subsection{Как появляется расстояние}
Когда в сети много частиц (\(N \gg 1\)), усреднение по большому числу путей \(\Gamma_{AB}\) позволяет сопоставить каждой паре координаты \(x^\mu_A\) и \(x^\mu_B\) так, что минимальное число посредников \(d_{AB}\) аппроксимируется геодезическим расстоянием в гладком многообразии:
\[
d_{AB} \approx \frac{1}{\ell_0} \int_{A}^{B} \sqrt{g_{\mu\nu} dx^\mu dx^\nu}
\]
где \(\ell_0\) — фундаментальная длина (планковский масштаб). Это приближение работает только для \textbf{больших расстояний} (\(d_{AB} \gg 1\)) и при наличии многих альтернативных путей.

Для малых \(d_{AB}\) или для изолированных пар частиц гладкое приближение не работает. Фундаментальная сеть не обязана подчиняться лоренцевой метрике.

\subsection{Световой конус как статистический эффект}
Лоренцева сигнатура и конечная скорость света \(c\) возникают из-за того, что в сети доминируют пути с большим числом посредников, передающих информацию последовательно. Для двух топологически идентичных частиц прямой канал через собственные моды \textbf{обходит} эту последовательную передачу.

Для наблюдателя, который воспринимает только эмерджентное пространство-время и не видит фундаментальных мод, передача через \(C_{AB}^{\text{(собств)}}\) выглядит как \textbf{мгновенная} (или со сверхсветовой скоростью). На самом деле она происходит вне причинной структуры эмерджентного пространства.

\subsection{Ограничения эмерджентной причинности}
В эволюционной версии сети вводится понятие \textbf{причинного горизонта} \(H_A(t)\): это множество частиц, достижимых из \(A\) по путям, длина которых не превышает \(D_{\max}(t)\). Величина \(D_{\max}(t)\) растёт со временем (моделируя расширение Вселенной).

Для топологически идентичных частиц прямая корреляция \(C_{AB}^{\text{(собств)}}\) существует \textbf{независимо от \(D_{\max}(t)\)}, потому что она не использует пути через посредников. Следовательно, такие частицы могут быть связаны даже если они находятся за пределами причинного горизонта в эмерджентном смысле.

\textbf{Это не нарушает причинность}, определённую в Теории Мод как \(\exists A \Leftrightarrow \exists B\): события \(A\) и \(B\) просто оказываются в одной замкнутой причинной петле, которая не обязана укладываться в рамки светового конуса.

\section*{Часть 3. Топологическая антенна: определение и свойства}
\addcontentsline{toc}{section}{Часть 3. Топологическая антенна: определение и свойства}

\subsection{Что такое антенна в Теории Мод}
Антенна — это устройство, состоящее из некоторого набора частиц (атомов, электронов проводимости), чья структура выделена так, что:
\begin{itemize}
    \item Она имеет внешние моды \(M_{\text{ант}\to j}\) для обмена энергией с излучением (в эмерджентном описании).
    \item Она обладает устойчивой собственной модой \(M_{\text{ант}\to \text{ант}}\) (петля, описывающая внутренние резонансы).
\end{itemize}
Обычная антенна (диполь, патч) опирается на эмерджентное пространство: её внешние моды возбуждают электромагнитные волны, распространяющиеся со скоростью \(c\).

\subsection{Топологическая антенна}
\textbf{Топологическая антенна} — это антенна, в которой:
\begin{enumerate}
    \item Собственная мода \(M_{A\to A}\) имеет нетривиальную топологию: например, чередование знаков поперечного спектра, соответствующее спину \(1/2\) (фермионная антенна) или спину \(1\) (бозонная).
    \item Материал антенны обладает \textbf{топологически защищёнными поверхностными состояниями} (как в топологических изоляторах Bi\(_2\)Se\(_3\), Bi\(_2\)Te\(_3\), (Sb,Bi)\(_2\)Te\(_3\) и других). Эти состояния устойчивы к примесям и дефектам.
    \item Поверхностные состояния формируют канал связи, в котором спин частицы жёстко связан с направлением в кристаллографическом пространстве (спин-моментное зацепление).
\end{enumerate}

В языке Теории Мод спин-моментное зацепление означает, что поперечный спектр \(F_{A\to A}(k)\) для поверхностной моды зависит от \textbf{направления в кристаллографических осях}, но не в эмерджентных. Для двух идентичных антенн, изготовленных из одного кристалла с одной ориентацией, их собственные моды когерентны, и эта когерентность сохраняется при любых перемещениях и поворотах антенн в эмерджентном пространстве.

\subsection{Требования к идентичности для двух режимов}

\paragraph{Режим 1: Передача энергии (без модуляции)}
Требуется:
\begin{itemize}
    \item Топологическая идентичность кристалла (пункты 1–3 раздела 1.5).
    \item Кристаллографическая идентичность (одинаковая ориентация решётки).
    \item Окружение может быть разным, температура и поля могут отличаться, но материал должен оставаться в топологической фазе.
\end{itemize}
При выполнении этих условий между антеннами существует постоянная корреляция \(C_{AB}^{\text{(собств)}}\), через которую может перетекать энергия при нарушении равновесия (например, если одна антенна подключена к источнику, а другая — к нагрузке). \textbf{Энергия передаётся независимо от расстояния и независимо от эмерджентной ориентации.}

\paragraph{Режим 2: Передача информации (модуляция корреляции)}
Требуется дополнительно:
\begin{itemize}
    \item \textbf{Идентичность ближнего электромагнитного окружения.} Наведённые токи, локальные поля, потенциалы соседних проводников, температура — всё это должно быть одинаково (или, по крайней мере, стационарно и взаимно известно) для обеих антенн. Это необходимо, чтобы модуляция собственной моды в передатчике не «гасилась» паразитными флуктуациями в приёмнике.
\end{itemize}
Практически это означает: две антенны из одной пластины помещаются в \textbf{идентичные экранированные камеры} с одинаковым питанием, одинаковой температурой и одинаковой разводкой сигнальных линий. Эмерджентное расстояние между камерами может быть любым — от метров до тысяч километров. Сами камеры могут быть ориентированы в пространстве как угодно.

\section*{Часть 4. Механизм сверхсветовой передачи}
\addcontentsline{toc}{section}{Часть 4. Механизм сверхсветовой передачи}

\subsection{Отсутствие задержки}
В классической электродинамике задержка сигнала между антеннами равна \(L/c\), где \(L\) — расстояние. В топологической связи \textbf{нет процесса распространения} в эмерджентном пространстве. Корреляция \(C_{AB}^{\text{(собств)}}\) существует всегда, независимо от \(L\).

Модуляция этой корреляции (например, изменение фазы или амплитуды собственной моды в антенне \(A\)) мгновенно отражается на корреляции, измеряемой в антенне \(B\). Единственная задержка — время перестройки собственной моды (время отклика антенны), которое не зависит от расстояния.

\textbf{Типичные времена отклика:} в топологических изоляторах — от терагерцового периода (\(\approx 10^{-13}\) с) до наносекунд (для медленных релаксационных процессов). Это может быть больше времени прохода света для \(L < 1\) м, но для \(L \gg 1\) м «сверхсветовая» передача очевидна.

\subsection{Передача энергии (без модуляции)}
Если антенна \(A\) подключена к источнику энергии (например, к выходу поляризатора вакуума), а антенна \(B\) — к нагрузке, то система стремится минимизировать глобальную нескомпенсированность. Поскольку корреляция между \(A\) и \(B\) максимальна, энергия перетекает из \(A\) в \(B\) без необходимости модуляции. Процесс похож на перетекание заряда между двумя одинаковыми конденсаторами, соединёнными идеальным проводником, но здесь роль проводника играет корреляция собственных мод.

\textbf{Скорость перетекания:} определяется импедансами нагрузки и источника, а не скоростью света. Для чисто активной нагрузки — экспоненциальный переход с постоянной времени \(\tau \sim R \cdot C_{\text{эффект}}\), где \(C_{\text{эффект}}\) — ёмкость корреляционного канала.

\subsection{Передача информации (модуляция)}
Для передачи информации в антенне \(A\) создаётся модуляция (амплитудная, фазовая или топологическая), которая изменяет \(C_{AB}^{\text{(собств)}}\). Приёмник \(B\) детектирует это изменение. Для успешного детектирования необходимо, чтобы флуктуации окружения (собственные шумы антенны \(B\)) были меньше уровня сигнала. Поэтому требуется идентичность окружения: если оба окружения одинаковы, то шумы синфазны и вычитаются при дифференциальном детектировании (или компенсируются).

\section*{Часть 5. Протоколы модуляции и детектирования}
\addcontentsline{toc}{section}{Часть 5. Протоколы модуляции и детектирования}

\subsection{Модуляция собственной моды (в передатчике)}
\begin{enumerate}
    \item \textbf{Фазовая модуляция}: изменение сдвига \(\delta\) в циклической последовательности \(F(k)\). Достигается подачей внешнего электрического поля, изменяющего спин-орбитальную связь в топологическом изоляторе. Скорость — до ТГц.
    \item \textbf{Амплитудная модуляция}: изменение величины \(|F(k)|\) (накачка неравновесными носителями заряда лазером). Скорость — ограничена временем рекомбинации носителей (пикосекунды–наносекунды).
    \item \textbf{Топологическая модуляция}: изменение индекса намотки собственной моды (рождение/исчезновение посредника в петле). Медленный процесс (предположительно, микросекунды), но очень помехоустойчив.
\end{enumerate}

\subsection{Детектирование (в приёмнике)}
Антенна \(B\) вычисляет корреляцию между своей собственной модой и сохранённым образом (или проводит гомодинное детектирование). Изменение \(C_{AB}^{\text{(собств)}}\) преобразуется в изменение напряжения или тока.

При идентичности окружения можно использовать дифференциальную схему: вычитать сигнал от холостой копии антенны (не модулированной), подавляя общие шумы.

\subsection{Синхронизация}
Поскольку сигнал распространяется без задержки, проблема синхронизации по времени отсутствует. Однако необходимо согласование \textbf{начальной фазы} собственных мод (циклический сдвиг \(\delta\)). Это решается процедурой «рукопожатия»: на короткое время антенна \(A\) подает фиксированную тестовую последовательность, антенна \(B\) подстраивает свой опорный генератор.

\section*{Часть 6. Экспериментальные предсказания}
\addcontentsline{toc}{section}{Часть 6. Экспериментальные предсказания}

\subsection{Независимость от эмерджентного расстояния (энергия)}
\textbf{Установка:} Две антенны из одной пластины Bi\(_2\)Te\(_3\), подключённые: \(A\) — к источнику постоянного напряжения (1 В, внутреннее сопротивление 50 Ом), \(B\) — к нагрузке 50 Ом. Расстояния: 0.1 м, 1 м, 10 м, 100 м.

\textbf{Измеряемая величина:} Мощность на нагрузке \(P_B\).

\textbf{Классика:} \(P_B \sim 1/L^2\) (если бы работала радиосвязь).

\textbf{Предсказание ТМ:} \(P_B\) не зависит от \(L\) (в пределах 5\%, ограниченных шумами). Задержка при включении/выключении не превышает времени отклика антенн (\(\sim 1\) нс) и не растёт с \(L\).

\subsection{Независимость от эмерджентной ориентации}
\textbf{Установка:} Антенны зафиксированы на расстоянии 1 м. Измеряется мощность на нагрузке. Затем одну антенну поворачивают в пространстве на 90°, 180°, 270° относительно другой.

\textbf{Предсказание:} Мощность не меняется (с точностью до погрешности). Это отличает топологическую связь от любого направленного радиоизлучения.

\subsection{Чувствительность к кристаллографической ориентации}
\textbf{Установка:} Две антенны, одна с ориентацией (111), вторая с (001) — вырезаны из разных пластин. Расстояние 1 м, одинаковое окружение.

\textbf{Предсказание:} Мощность близка к нулю, даже если антенны эмерджентно сориентированы «правильно». Это доказывает, что связь не электромагнитная.

\subsection{Зависимость от идентичности окружения (модуляция)}
\textbf{Установка:} Две антенны из одной пластины, расстояние 1 м. Передатчик модулирует сигнал (например, амплитудная модуляция 1 кГц). Приёмник синхронно детектирует. Окружение первой антенны — стабильно. Окружение второй антенны варьируется: (а) то же самое (температура 25°C, экранирование); (б) другое (температура 40°C, рядом работающий двигатель).

\textbf{Предсказание:} В случае (а) коэффициент ошибок (BER) низкий, отношение сигнал/шум > 20 дБ. В случае (б) BER высокий, связь возможна только при снижении скорости (из-за увеличения шумов). Это показывает, что передача информации требует идентичности окружения, а передача энергии — нет.

\subsection{Передача энергии на расстояние без проводов (демонстрация)}
\textbf{Установка:} Поляризатор вакуума (фрактальная решётка) на 1 Вт подключён к антенне \(A\). Антенна \(B\) находится в другой комнате (или в другом здании) на расстоянии 100 м и подключена к светодиоду (3 В, 20 мА). Обе антенны вырезаны из одной пластины.

\textbf{Предсказание:} Светодиод горит полным накалом, яркость не зависит от расстояния (при условии, что антенны изготовлены идентично). Задержка включения после включения поляризатора — менее 1 мкс.

\section*{Часть 7. Применения}
\addcontentsline{toc}{section}{Часть 7. Применения}

\subsection{Глобальная энергетическая сеть без проводов}
Массовое производство антенн из одной пластины (тысячи штук) позволяет снабжать энергией удалённые регионы без прокладки кабелей и без потерь в линиях. Каждая антенна напрямую коррелирует с любой другой из той же партии. Достаточно иметь один мощный поляризатор вакуума — и все устройства с антеннами получают энергию.

\subsection{Защищённая связь}
Поскольку для информационной связи требуется не только кристаллографическая идентичность, но и идентичность ближнего окружения (которую трудно скопировать без доступа к физической установке), возможна физическая защита канала. Даже если злоумышленник имеет антенну из той же пластины, но не может воспроизвести окружение (температуру, экранирование, разводку), он не сможет демодулировать сигнал без внесения ошибок.

\subsection{Связь с подводными и подземными объектами}
Поскольку сигнал не распространяется через эмерджентное пространство, он не затухает в воде, земле, металлах. Связь с подводными лодками, шахтами, бункерами становится такой же простой, как связь в лаборатории. Единственное требование — идентичность антенн и одинаковая разводка окружения.

\subsection{Распределённые вычисления без задержек}
Если процессорные ядра оснастить топологическими антеннами, они могут обмениваться данными без задержки на расстояние (в пределах одной пластины — естественно, на разных пластинах — только если они вырезаны из одной заготовки). Это может радикально изменить архитектуру суперкомпьютеров.

\subsection{Ограничения}
\begin{itemize}
    \item Требование одинаковой кристаллографической ориентации ограничивает размер антенн одной пластиной (максимум 300 мм в диаметре для кремниевой технологии, для Bi\(_2\)Te\(_3\) — меньше). Для глобальной сети нужно либо разрезать пластину и развозить кусочки по всей планете (что возможно), либо научиться воспроизводить кристаллографическую ориентацию между разными пластинами с точностью до 0.01°, что пока нереализуемо.
    \item Для передачи информации требуется идентичность окружения, что ограничивает мобильность. Приёмник должен находиться в стабильных условиях, эквивалентных условиям передатчика. Для стационарной инфраструктуры это выполнимо.
\end{itemize}

\section*{Часть 8. Связь с фундаментальными принципами Теории Мод}
\addcontentsline{toc}{section}{Часть 8. Связь с фундаментальными принципами Теории Мод}

\subsection{Двусторонняя причинность}
Топологическая связь является прямым проявлением Аксиомы 13: для любой причинно связанной пары \((A, B)\) существует замкнутый путь с \(\sum F = 0\). Идентичные антенны гарантируют, что этот путь существует всегда, независимо от эмерджентного расстояния.

\subsection{Компенсация как аттрактор}
Передача энергии через топологическую связь — это релаксация глобальной нескомпенсированности. Когда одна антенна получает энергию от внешнего источника, система перераспределяет её на другие идентичные антенны, возвращаясь к равновесию.

\subsection{Нелокальность}
Зависимость корреляции от расстояния отсутствует, потому что фундаментальная сеть мод не имеет метрики в том смысле, в каком её имеет эмерджентное пространство. Это не нарушает локальность, а скорее показывает, что «локальность» — эмерджентное свойство, не применимое к топологически идентичным объектам.

\section*{Часть 9. Заключение}
\addcontentsline{toc}{section}{Часть 9. Заключение}

Топологическая связь — прямое предсказание Теории Мод, которое может быть проверено в ближайшие годы. Для двух идентичных топологических антенн, изготовленных из одного кристалла топологического изолятора с одинаковой кристаллографической ориентацией, существует прямой корреляционный канал, который:
\begin{itemize}
    \item Не зависит от эмерджентного расстояния (сотни метров — тысячи километров).
    \item Не зависит от эмерджентной ориентации (антенны могут быть повёрнуты как угодно).
    \item Позволяет передавать энергию без модуляции, даже при разном окружении.
    \item Позволяет передавать информацию только при идентичности ближнего электромагнитного окружения (температура, наведённые токи, разводка).
\end{itemize}

Эффект может быть использован для создания глобальной беспроводной энергетической сети, где один поляризатор вакуума питает миллионы устройств. Основное технологическое ограничение — необходимость вырезать все антенны из одной пластины, чтобы гарантировать кристаллографическую идентичность. Для стационарной инфраструктуры это выполнимо: пластина режется на чипы, чипы разносятся по стране, каждый подключается к своему устройству.

Документ самодостаточен. Экспериментальные предсказания (раздел 6) легко воспроизводимы в лаборатории, имеющей доступ к технологии роста тонких плёнок топологических изоляторов и электронной литографии. Теория Мод будет подтверждена или опровергнута результатами этих экспериментов.

\end{document}